%% =========================================================================
\section{Variants of the Conjecture}
\label{sec:variants}
%% =========================================================================

The Euclid--Mullin construction makes sense with any factor-selection
rule: at each step, $\Prod{2}(n)+1$ has a non-empty set of prime
factors, and the rule picks one.  Mullin's Conjecture is specific to
the $\minFac$ rule.  This section collects the known variants---some
easy, some provably incomplete, one almost closed---and uses them to
locate standard MC within a broader landscape.

\begin{roadmap}
\begin{enumerate}[nosep,leftmargin=*]
\item[\S\ref{sec:target-ladder}] The target ladder --- how little must be proved; down to (V), the weakest orbit statement.
\item[\S\ref{sec:composite-floor}] The smallness family $\mathrm{MissingFinite} \Rightarrow \mathrm{WM} \Rightarrow \mathrm{RD}$, and the elementary \emph{floor} beneath all of it.  That floor, $(\mathrm{C}\infty)$, is the subject of \S\ref{sec:cinfty}.
\item[\S\ref{sec:random-variant}] Random-factor variant --- MC by mixing; independence vs.\ selection trade-off.
\item[\S\ref{sec:factor-tree}] Factor tree and GenMixedMC --- nondeterministic access to all prime factors.
\item[\S\ref{sec:almost-all-mixed}] Almost all GenMixedMC --- per-prime density-$1$ result plus diagonal, \emph{unconditional} (the Dirichlet-density input is itself machine-checked).
\item[\S\ref{sec:minfac-secondminfac}] $\{\minFac,\operatorname{secondMinFac}\}$ variant --- two-point spectral contraction, UFDStrong.
\item[\S\ref{sec:ensemble-variant}] Ensemble MC --- almost every starting point, backward dynamics.
\item[\S\ref{sec:weak-mullin}] Weak Mullin --- analytic weakenings, partial results.
\item[\S\ref{sec:function-field}] Function field analog --- $\FF_p[t]$, Weil bound, orbit-specificity barrier; \emph{false} over $\FF_5[t]$ (stable Sylvester tower).
\end{enumerate}
\end{roadmap}

\paragraph{The Selection Rule Landscape}
\label{sec:selection-landscape}

Fix an initial prime $p_0 = 2$.  Given a \emph{selection rule}
$\sigma$ that, for any integer $N \geq 2$, chooses a prime factor
$\sigma(N) \mid N$, define:
\begin{align}
  P_\sigma(0) &= 2, &
  P_\sigma(n{+}1) &= P_\sigma(n) \cdot \sigma(P_\sigma(n)+1).
\end{align}
The selection rule determines the sequence; the running product
$P_\sigma(n)$ determines the sequence history.
Five rules are studied in the formalization:

\smallskip
\begin{center}
\begin{tabular}{@{}l@{\quad}l@{\quad}l@{}}
\toprule
\textbf{Selection rule} $\sigma$ & \textbf{MC status} & \textbf{Key property} \\
\midrule
Random (uniform) & Open (heuristically easy) & Randomised selection \\
$\minFac$ (standard EM) & \textbf{Open} (CME) & Deterministic, lethal \\
$\{\minFac, \operatorname{secondMinFac}\}$ & \textbf{Open} (UFDStrong) & Two-point diversity \\
$(1{-}\varepsilon)\,\minFac + \varepsilon\,$random & a.a.\ (unconditional) & Full factor access \\
Steered (Booker~\cite{BoSa2012}) & True (by construction) & Adversarial \\
$\mathrm{maxFac}$ (Cox--vdP~\cite{CoxVdP1968}) & False & Misses primes \\
\emph{every} factor (Sylvester) & \textbf{False} & Autonomous closure \\
\bottomrule
\end{tabular}
\end{center}

\smallskip\noindent
The landscape reveals a tension: the random rule is the easiest to
analyze (mixing theory) but the worst at selection (probability
$1/|F|$ per hit); $\minFac$ is the best at selection (lethal hits)
but the hardest to analyze (deterministic, no independence); the
$\{\minFac, \operatorname{secondMinFac}\}$ variant sits between them,
offering a two-point spectral contraction that is provably sufficient
under a single open hypothesis; and the factor tree framework
(\S\ref{sec:factor-tree}) generalizes all three by allowing access
to the full factor set at every step, reducing GenMixedMC to sieve
theory for almost all starting points.

\paragraph{Growth against capture: the take-all rule.}
It is tempting to think that a \emph{more generous} rule---one that
banks every prime factor of the Euclid number rather than a single
one---could only help, and that if even that rule missed primes then
$\minFac$ would have no chance.  The opposite is true, and the
computation is short enough to record.

Banking every factor gives
$P_{n+1} = P_n \cdot \operatorname{rad}(P_n+1)$, and whenever the
Euclid number is squarefree---as it is for every computed term---this
is
\[
  P_{n+1} \;=\; P_n\,(P_n+1),
\]
i.e.\ $2, 6, 42, 1806, 3263442, \dots$: \emph{Sylvester's sequence},
shifted by one.  The rule has closed the recursion into an autonomous
polynomial map, and modulo~$q$ the walk becomes
$w \mapsto w^2 + w$ with no dependence on anything but~$w$.  Death
requires $w^2 + w = -1$, that is $w^2+w+1 = 0$, whose roots are the
primitive cube roots of unity---present in $\FF_q$ exactly when
$q \equiv 1 \pmod 3$.  So the take-all rule misses \emph{every} prime
$q \equiv 2 \pmod 3$ \emph{as long as its Euclid numbers stay
squarefree}---an assumption (whether every Sylvester number is squarefree
is open~\cite{Vardi1991}) that we do not formalize; the statement is
therefore conditional, not a theorem of the development.

The reason is growth.  Ordering the rules by how fast they enlarge the
accumulator,
\[
  \minFac \;<\; \cdots \;<\; \mathrm{maxFac} \;<\; \text{take-all},
\]
and ordering them by how badly they fail gives the \emph{same} order.
Each prime~$q$ needs many steps at which $q$ might divide $P_n+1$;
the faster the accumulator grows, the fewer such trials it gets.  This
is the greedy-for-slow-growth and $0$-accessibility analysis of
\S\ref{sec:bag} seen from the far end: $\minFac$ is not merely one rule
among several, it is the extreme point of the only axis that matters,
and the take-all rule is the opposite extreme where the dynamics
degenerates into an algebraic obstruction.

\smallskip\noindent
This also isolates what generosity was worth having.  The take-all rule
buys the removal of the \emph{selection} barrier and pays for it with
\emph{growth}.  The next subsection keeps the first and refuses the
second.

\smallskip\noindent
The map $w \mapsto w^2+w$ found here is not confined to this variant.  It is
also the residue walk of the $\minFac$ rule on the perpetual-primality
branch (\S\ref{sec:autonomous-branch}) and, conjugated by $w \mapsto w+1$,
the map governing $(\mathrm{C}\infty)$ itself (\S\ref{sec:level-tower}).

%% =========================================================================
\subsection{The Target Ladder: How Little Must Be Proved?}
\label{sec:target-ladder}
%% =========================================================================

Keep the $\minFac$ accumulator---so the growth advantage is
untouched---but count a prime as captured as soon as it \emph{divides}
a Euclid number, whether or not it is the factor selected.  This is a
weakening of the \emph{question}, not a change of rule, and it yields

\begin{definition}[\defname{(V)} ---
\lean{EM/Equidist/WeakHitting.lean}{EveryPrimeDividesEuclid}]
\label{def:V}
Every odd prime divides some Euclid number: for each prime $q \neq 2$
there is an~$n$ with $q \mid \Prod{2}(n)+1$; equivalently
$\mathrm{HittingSet}(q) \neq \emptyset$.  The parity restriction is
forced---$\Prod{2}(n)$ is always even, so no Euclid number is even
and~$2$ is the seed rather than a capture.
\end{definition}

\noindent
(V) slots into a chain of successively weaker targets, all of them
statements about the \emph{same} orbit:

\smallskip
\begin{center}
\begin{tabular}{@{}l@{\quad}l@{}}
\toprule
\textbf{Target} & \textbf{What it asks of the walk mod~$q$} \\
\midrule
$\mathrm{HH}$ & hits $-1$ \emph{cofinally often} \\
$\mathrm{SingleHit}$ & hits $-1$ \emph{once past the sieve gap} \\
$\MC$ & is captured, i.e.\ hits $-1$ where $q$ is minimal \\
$(\mathrm{V})$ & hits $-1$ \emph{at least once, ever} \\
$\mathrm{GenMixedMC}(2)$ & hits $-1$ somewhere in the factor tree \\
\bottomrule
\end{tabular}
\end{center}

\smallskip\noindent
The implications run downwards and are machine-verified
(\lean{EM/Equidist/WeakHitting.lean}{weak\_hitting\_ladder}).  The
repository already carried the outer rungs; (V) was the missing one.
Matching the ladder of targets is a ladder of hypotheses: the
one-horizon Fourier criterion of \S\ref{sec:hard} asks for a good
window past \emph{every} stage because $\MC$ needs the
first-missing-prime bootstrap, whereas (V) follows from a single window
per prime, anywhere
(\lean{EM/Equidist/WeakHitting.lean}{oneWindowGain\_implies\_V}).

\begin{aside}
\textbf{How much weaker is (V), really?}  Less than it looks.  Suppose
every prime below~$q$ has appeared by stage~$N_0$.  Then for $n \geq
N_0$ no prime below~$q$ divides $\Prod{2}(n)+1$, so a hit at such an~$n$
forces $\minFac(\Prod{2}(n)+1) = q$ and $q$ \emph{is} captured.  So
$\MC$ and (V) differ only on hits occurring \emph{before} the sieve
gap---at most $\pi(q)$ steps, since each such hit is shielded by a
distinct smaller prime
(\lean{EM/Population/HittingSetStructure.lean}{hittingSet\_ncard\_le\_appearing}).
Widening the notion of capture buys a finite prefix per prime, not a
different problem: the bottleneck was never selection, which is free
past the gap, but hitting.
\end{aside}

%% =========================================================================
\subsection{The Smallness Family and Its Floor}
\label{sec:composite-floor}
%% =========================================================================

Below $\MC$ sits a second family of weakenings, orthogonal to the target
ladder of \S\ref{sec:target-ladder}.  Instead of asking less of the
walk, they ask only that the set of \emph{missed} primes
$\mathcal{M} = \{q \text{ prime} : \forall k,\ \seq{2}(k) \neq q\}$ be small:

\begin{center}
\begin{tabular}{@{}l@{\quad}l@{}}
\toprule
\textbf{Statement} & \textbf{Assertion} \\
\midrule
$\mathrm{MissingFinite}$ & $\mathcal{M}$ is finite \\
$\mathrm{WM}$ (Weak Mullin) & $\sum_{q \in \mathcal{M}} 1/q$ converges \\
$\mathrm{RD}$ & $\sum_k 1/\seq{2}(k)$ diverges \\
\bottomrule
\end{tabular}
\end{center}

\smallskip\noindent
These are ordered: $\MC \Rightarrow \mathrm{MissingFinite} \Rightarrow
\mathrm{WM} \Rightarrow \mathrm{RD}$, the last step because the primes
split into those that appear and those that do not, and
$\sum_p 1/p$ diverges
(\lean{EM/Population/WeakMullin.lean}{wm\_implies\_rd}).  Every one of
these implications runs \emph{downwards} from $\MC$.  Until now the
family had no lower bound at all: nothing in it was known
unconditionally, and nothing was known to require anything.

\paragraph{The floor.}
It has one, and the argument is elementary.  Suppose $\Prod{2}(n)+1$ is
prime for every $n \geq N$ (\emph{perpetual primality}).  Then the least
factor of the Euclid number is the Euclid number, so
$\seq{2}(n+1) = \Prod{2}(n)+1$, and the accumulator's growth---already
geometric for trivial reasons, since every term is at least $2$---is
inherited by the selected primes:
\[
  2^{\,n+1} \leq \Prod{2}(n)
  \quad\text{unconditionally}
  \qquad\Longrightarrow\qquad
  \sum_k \frac{1}{\seq{2}(k)} < \infty
  \ \text{ on this branch.}
\]

\begin{keyresult}
\textbf{The composite floor}
(\lean{EM/Population/CompositeFloor.lean}{composite\_floor\_landscape}).
Write \textbf{(C$\infty$)} for the statement that $\Prod{2}(n)+1$ is
composite for infinitely many~$n$.  Then
\[
  \mathrm{RD} \Rightarrow (\mathrm{C}\infty), \quad
  \mathrm{WM} \Rightarrow (\mathrm{C}\infty), \quad
  \mathrm{MissingFinite} \Rightarrow (\mathrm{C}\infty), \quad
  \MC \Rightarrow (\mathrm{C}\infty),
\]
and on the perpetual-primality branch every one of the three smallness
statements fails outright.  \textbf{(C$\infty$) is open.}
\end{keyresult}

\noindent
Note how little is consumed.  The autonomous-branch analysis of
\S\ref{sec:hard} derives $\lnot\MC$ on this branch through Bertrand's
postulate, and a density-$1/2$ failure through $\Phi_3$ having no root
modulo $q \equiv 2 \pmod 3$.  Neither is needed here: geometric growth of
the accumulator alone kills the reciprocal-sum statements, which sit far
below $\MC$.  So $\MC \Rightarrow (\mathrm{C}\infty)$ is recovered from a
strictly weaker hypothesis by a strictly more elementary route.

\smallskip\noindent
That floor is the subject of \S\ref{sec:cinfty}, which is where the rest of
this paper's unconditional work on $(\mathrm{C}\infty)$ lives.  Two remarks
belong here, with the smallness family itself.

\smallskip\noindent
One unconditional inequality survives in the other direction.  Every
prime below the least missed prime \emph{has} been selected, so it
contributes to the sequence's reciprocal sum:
\[
  \sum_{p < \min \mathcal{M}} \frac{1}{p}
  \;\leq\; \sum_k \frac{1}{\seq{2}(k)}
  \qquad
  \text{(\lean{EM/Population/CompositeFloor.lean}{sum\_inv\_primes\_below\_le\_tsum})}
\]
whenever the right-hand side converges.  It runs opposite to
$\mathrm{WM} \Rightarrow \mathrm{RD}$: there, thinness of the missed set
forces the sum to diverge; here, a convergent sum caps how far the first
gap can be.


\smallskip\noindent
And what the family actually needs is stronger than $(\mathrm{C}\infty)$:
not that the Euclid numbers are sometimes composite, but that they are
sometimes \emph{smooth relative to their own size}
(\S\ref{sec:growth-floor}).

%% =========================================================================
\subsection{The Random-Factor Variant}
\label{sec:random-variant}

Consider a variant where, instead of the smallest prime factor, one
picks a \emph{random} prime factor of~$\Prod{2}(n)+1$ at each step.
Heuristically MC should be easy here, and one is tempted to say ``by
classical mixing theory'': the residue $r_n = \Prod{2}(n) \bmod p$
performs a multiplicative walk on $(\ZZ/p\ZZ)^\times$ whose multipliers
are random elements of the group.  But the steps are \emph{not}
independent and \emph{not} independent of the position: the multiplier is
a uniformly random prime factor of the specific integer
$\Prod{2}(n)+1$ (if that integer is prime the step is forced), so
Diaconis--Shahshahani does not apply as stated (Dead End~\#47).  What is
formalized is the nondeterministic content only: purification,
$\mathrm{PureRandomMC} \leftrightarrow \mathrm{MixedMC}$
(\code{EM/Stochastic/RandomFactorMC.lean}); the probabilistic statement
``$p$ is eventually selected almost surely'' is \textbf{not} formalized
and we do not claim it.

\paragraph{Comparison with $\minFac$.}
The actual EM sequence uses
$\seq{2}(n{+}1) = \minFac(\Prod{2}(n){+}1)$, a \emph{deterministic}
function of the walk position.  Compared to the random-factor variant,
$\minFac$ is strictly better in every measurable dimension:

\smallskip
\begin{center}
\begin{tabular}{l@{\quad}l@{\quad}l}
\toprule
 & \textbf{Random factor} & $\boldsymbol{\minFac}$ \\
\midrule
\textbf{Selection when $q \mid P(n){+}1$}
  & Probability $1/|F(S_n)|$ & Probability~1 \\
\textbf{Accumulator growth}
  & Faster (random $\geq \minFac$) & Slowest possible \\
\textbf{Hits on $-1$ needed}
  & Infinitely many & \textbf{One} (lethal hit) \\
\bottomrule
\end{tabular}
\end{center}

\smallskip\noindent
\emph{Selection efficiency.}\enspace Past the sieve gap, 0-accessibility
guarantees: if $q \mid \Prod{2}(n)+1$ then
$q = \minFac(\Prod{2}(n)+1)$---selection is certain, not
probabilistic.  \emph{Growth rate.}\enspace Since
$\minFac(N) \leq p$ for any factor~$p$ of~$N$, the accumulator
$\Prod{2}(n)$ grows minimally under the $\minFac$ rule, giving
exponentially more steps before it becomes astronomically large---more
chances for~$q$ to divide.  \emph{Lethality.}\enspace The inductive
bootstrap shows one single hit suffices: $\MC({<}\,q)$ ensures~$q$
is minimal whenever it divides, so hitting $-1$ once means~$q$
enters the sequence.  The random variant requires infinitely many
hits to overcome low per-hit selection probability.

\emph{MinFac is strictly better in every measurable dimension.
So why can't we prove it works?}

\paragraph{The independence trade-off.}
The answer is surgical: the random variant possesses a single property
that $\minFac$ lacks, and that property is the \emph{only thing the
proof uses}.

The property is \textbf{independence of multiplier from walk
position}.  In the random-factor variant, the multiplier at step~$n$
is selected independently of which residue class
$\Prod{2}(n) \bmod q$ occupies.  Walk position and multiplier are
independent random variables, and independence is the engine that
drives mixing.  Diaconis--Shahshahani's theorem applies whenever the
multipliers are i.i.d.\ from a distribution whose support generates
the group.

In the $\minFac$ variant, the multiplier is
$\minFac(\Prod{2}(n){+}1) \bmod q$---a deterministic function of the
full integer~$\Prod{2}(n)+1$.  The CRT multiplier invariance
(\lean{EM/Group/CRT.lean}{crt\_multiplier\_invariance}) shows
that the outcome is independent of $\Prod{2}(n) \bmod q$ at the
\emph{population} level: for generic integers, knowing the mod-$q$
residue does not predict $\minFac$.  But for the specific orbit~$n=2$,
the multiplier is a deterministic value, and mixing theorems for
deterministic sequences cannot be invoked.

The random variant thus sacrifices selection efficiency (it may
miss~$q$ even when $q$ divides) in exchange for statistical
independence.  The $\minFac$ variant has perfect selection but no
independence:

\smallskip
\begin{center}
\begin{tabular}{l@{\quad}l@{\quad}l}
\toprule
 & \textbf{Random factor} & $\boldsymbol{\minFac}$ \\
\midrule
Selection & Bad ($1/k$ chance) & Perfect \\
Mixing & Random walk theorem & Unproved \\
\bottomrule
\end{tabular}
\end{center}

\smallskip\noindent
The random variant's proof: mixing $\to$ cofinal hitting $\to$
infinitely many chances compensate for low per-hit probability.
The $\minFac$ variant needs only one hit---absurdly less than the
random variant achieves---but has no way to produce even that one
without mixing.

\paragraph{The stochastic domination failure.}
\label{sec:stoch-dom}
The natural instinct is to argue by \emph{stochastic domination}:
if the random walk visits~$-1$ cofinally, and $\minFac$ is ``better
in every way,'' then $\minFac$ should also visit~$-1$.  This argument
fails because the two processes \emph{diverge after one step}.

At step~$n$, both processes face the same Euclid number
$\Prod{2}(n)+1$ with factor set~$F(S_n)$.  The random variant picks a
uniformly random element of~$F(S_n)$; the $\minFac$ variant picks the
minimum.  After this step, the two processes have different
accumulators (they multiplied by different primes), so they face
different Euclid numbers at the next step.  The processes explore
completely different trajectories.  Domination fails because the
two walks live on different probability spaces after step~1; we
cannot compare them step-by-step.

A subtler comparison uses \emph{visit measures}.  Define
$\mu_{\mathrm{rand}}(a) = \lim_{N} \frac{1}{N}|\{n < N :
w_q^{\mathrm{rand}}(n) = a\}|$ and analogously
$\mu_{\min}(a)$ for the $\minFac$ walk.  For the random variant,
$\mu_{\mathrm{rand}} = $ uniform on $(\ZZ/q\ZZ)^\times$ by mixing.
For $\minFac$, CME/CCSB asserts $\mu_{\min} \approx$ uniform.  The
comparison $\mu_{\min}(-1) \geq 1/(q{-}1)$ is precisely what needs
proof, and is the content of CME\@.

